The determinant identity uses the actual diamond formula, not merely
the reduction of its eigenvalues. For a good prime $r$,
$d(\operatorname{Fr}_r)=r\langle r\rangle$. The formula above shows
that $(d(\operatorname{Fr}_r)-1)F_U$ reduces to a sum of integral
multiples of the undivided corrections $\theta B_U[J]\Theta_J$:
the scalar $r\psi_U(r)-1$ lies in $\mathfrak m_{\rm coef}$.
A fresh trace kills every such correction. Indeed, to test at $n$
choose $q\nmid nJ$; then $nq$ cannot have squarefree part $J$, and
$\theta\equiv1\pmod{\mathfrak m_{\rm coef}}$. This remains true
after further fresh traces. Continuity and the closed integral
module therefore give, on the reduced trace orbit of $F_U$,
\[
 d(u)t(v)=t(v),\qquad
 t(u)t(v)=t(uv)+t(uv^{-1}),\qquad t(1)=0.
\]
The middle equality follows from the characteristic-zero matrix
identity
$\operatorname{tr}(A)\operatorname{tr}(B)
=\operatorname{tr}(AB)+\det(B)\operatorname{tr}(AB^{-1})$,
using the determinant identity just proved; the last follows from
$\operatorname{tr}(1)=2$.

Two fresh primes outside $S\cup U$ contribute two to the squarefree
weight in (O.9). Thus $B_U[nq_1q_2]$ is divisible by four, no unary
correction occurs there, and $H_U[nq_1q_2]$ remains divisible by two.
Here $R\ge2$ and the actual-stage budget were fixed before the test.
It follows that $t(g)t(h)F_U=0$ after reduction. Applying the middle
identity to $t(gh)t(h)$ gives $t(gh^2)F_U=t(g)F_U$.
Invariance under all squares, and continuity, give factorization
through the elementary quadratic quotient unramified outside
$S\cup U$. Over $\Q$ this quotient is finite: its squareclasses
are generated by $-1$, $2$, and the primes in that support.
This does not assert abelianity of the characteristic-zero
Galois representations.

At $p\in U$, first keep the summands of $F_U$ separately.
The block $\theta B_U$ and each unary block with $p\in J$ have
level exponent at most one and primitive quadratic character at $p$.
The unary blocks with $p\notin J$ are degeneracies of forms at a
level omitting $p$. This distinction is preserved when retaining all
old multiplicities. In particular, the latter blocks cannot acquire a
Steinberg constituent merely by being viewed at a larger level.

Enlarge the characteristic-zero coefficient field to realize all the
cuspidal newform constituents and the Eisenstein character pairs of
these blocks. On a constituent and all its degeneracies, define $t(g)$
and $d(g)$ to be the scalar trace and determinant of its two-dimensional
Galois representation, acting as the identity on its multiplicity
space. Distinct characteristic-zero constituents with the same good
eigenvalues have the same trace and determinant by Chebotarev, so
these definitions agree on coinciding summands. In particular they
give well-defined operators on their sum. No decomposition modulo
two, and no residual irreducibility assertion, enters this construction.

